\documentclass[a4paper]{article}
\input{Packages.tex}
\begin{document}
\title{Technical Note: \\
\emph{Social Security Design and Its Political Support}}

\title{Technical Note: \\
\emph{Social Security Design and Its Political Support}}

\author{Rasmus K.\ Berg\thanks{University of Copenhagen. Oster Farimagsgade 5, 1353 Copenhagen K, Denmark. E-mail: \texttt{rasmus.kehlet.berg@econ.ku.dk}.}
\and
Mart\'{\i}n Gonzalez-Eiras\thanks{University of Bologna. Piazza Scaravilli 2, 40126 Bologna. E-mail: \texttt{mge@alum.mit.edu}. Web: \texttt{alum.mit.edu/www/mge}.}}
\maketitle

\begin{abstract}
\noindent
The note formally presents a number of different models used in the paper \textit{Social Security Design and Its Political Support}. In many instances, the note is more general than in the final paper. 

Each model is documented in two stages: The first one describes and characterizes both the economic equilibrium and the politico-economic equilibrium. In the second stage, we go through more carefully how we solve the model numerically, including the solution algorithms we apply throughout. All model documentations are self-contained, with no cross-references. This means that a lot of the text will naturally be identical across model types.

At the current stage, models include the following overall variations: The analytical-informal model, the informal savings model, and the US model. For each of the three, we currently implement a number of different solutions. First, we have one set of implementations with one common $X$ across all types and one with type-specific ones, $X_j$. Second, we have one where pension characteristics are fixed (exogenous), one where they are chosen permanently in one period, and one where characteristics are chosen without commitment equivalent to the tax rate.
\end{abstract}


\listofmodels

% First model
\subimport{informalAnalytical/}{main}

% Reset numbering:
\resetdocumentnumbering
\subimport{informalSavings/}{main}

% Reset numbering:
\resetdocumentnumbering
\subimport{US/}{main}



\end{document}